\chapter{Registro de uso de inteligencia artificial generativa}

\label{anx:ia}

En este anexo se indican ejemplos agrupados de los usos de inteligencia artificial generativa realizados durante la elaboración de este trabajo, indicando la herramienta usada, el objetivo de su uso, el prompt empleado junto con los parámetros de entrada y el resultado obtenido, en caso de que sea necesario.
No se incluyen como resultados obtenidos aquellos que ya están incluidos en la memoria, sino que se hace una referencia a la sección de la memoria correspondiente.
 
\begin{fichaia}{FichaIA. Registro de trazabilidad y generación de metadatos de IA}
  \fichadato{Fecha}{30/09/2026}
  \fichadato{Herramienta}{Gemini (Google)~\parencite{gemini}}
  \fichadato{Fase}{WP1. Definición y objetivos}
  \fichadato{Objetivo}{Ofrecer de forma transparente la información sobre el uso de inteligencia artificial generativa utilizada por el autor, incluyendo los prompts y ajustes realizados.}
  \fichadato{Resultado}{Código de ficha en LaTeX (\texttt{\textbackslash begin\{fichaia\}}) normalizado que compila los comandos, parámetros y modificaciones de la sesión.}
  \fichadato{Validación}{Elaboración propia y verificación de la estructura exigida para el anexo de transparencia de la memoria del TFM.}
  \fichaprompt{A partir de la conversación o de las modificaciones que acabamos de realizar, genera la ficha de registro en formato LaTeX \textbackslash begin\{fichaia\}... (incluyendo la ficha completa del componente en Latex)}
\end{fichaia}

\begin{fichaia}{Sistema de diseño de la memoria}
  \fichadato{Fecha}{29/09/2026 -- 18/10/2026}
  \fichadato{Herramienta}{Claude (Anthropic)~\parencite{claude}}
  \fichadato{Fase}{WP1. Definición y objetivos}
  \fichadato{Objetivo}{Disponer de un sistema de diseño para la memoria, a partir de los estilos de la universidad, que dé una visualización homogénea y facilite generar informes, gráficas, infografías, tablas, mapas y, después, árboles de decisión y configuraciones de XGBoost.}
  \fichadato{Resultado}{Guía de estilo, \texttt{tfm-estilo.sty} (v1.8) y listado de componentes de \LaTeX{}: colores accesibles, castellano, dos columnas, diagramas, Kanban, gráficos, citas, notas, comparación real frente a predicho, estilo \texttt{tfmarbol}, \texttt{\textbackslash esquemaboosting} y módulo \texttt{arboles.py}, que genera árboles a partir de modelos de scikit-learn y XGBoost.}
  \fichadato{Validación}{Compilación con la plantilla base de la universidad; colores tomados de tokens verificados (WCAG 2.1 y simulación de daltonismo); prueba de \texttt{arboles.py} con modelos de scikit-learn y XGBoost entrenados sobre datos sintéticos; revisión visual del autor y comprobación a medida que se incorporaban los componentes a la memoria.}
  \fichaprompt{Estoy haciendo un trabajo fin de master del máster de ciencia de datos de la UOC, necesito un sistema de diseño para poder generar informes, gráficos, infografías, presentación de resultados, diagramas de arquitectura de redes neuronales, arquitectura y diseño de código. La temática del proyecto es la utilización de imágenes aéreas de alta resolución para extraer información de las calles mediante deep learning, segmentación semántica y visión por computador, para permitir clasificar las calles en base a su accesibilidad, para poder ayudar a generar los planes de accesibilidad con información sacda de fotos, y reducir el gasto y tiempo que lleva la inspección in situ. Te paso unas capturas del estilo definido por la universidad para la entrega de la memoria. Evidentemente el formato de la memoria es académico, pero podría llegar a tener algún cambio si así lo solicito. ¿Me puedes crear el sistema de diseño? Si tienes alguna pregunta previa, me dices}
  \fichaprompt{Necesito que actualices el system design para incluir un componente que permita pintar árboles de decisión, así cómo configuraciones de XG Boosting}
\end{fichaia}

\begin{fichaia}{Figuras y tablas de introducción y planificación}
  \fichadato{Fecha}{30/09/2026 -- 03/10/2026}
  \fichadato{Herramienta}{Claude (Anthropic)~\parencite{claude}; Gemini (Google)~\parencite{gemini} para la figura~\ref{fig:proceso_actual_propuesto}}
  \fichadato{Fase}{WP1. Definición y objetivos}
  \fichadato{Objetivo}{Mejorar la comprensión de la introducción con visualizaciones, y adaptar al \emph{design system} el diagrama de Gantt y la tabla de planificación elaborados antes de forma manual por el autor, sin modificar su contenido.}
  \fichadato{Resultado}{Código TikZ de la infografía de tres cifras (figura~\ref{fig:infografia_accesibilidad}), del proceso actual y propuesto (figura~\ref{fig:proceso_actual_propuesto}), del ciclo de la metodología (figura~\ref{fig:metodologia_ciclo}, con la etapa de desarrollo resaltada después por el autor) y de los objetivos MVP (figura~\ref{fig:objetivos_mvp}); código del Gantt (figura~\ref{fig:gantt_tfm}) y de la tabla de planificación (tabla~\ref{tab:planificacion_tfm}). Ajustes menores de redacción.}
  \fichadato{Validación}{Infografía: cifras aportadas por el autor a partir del INE (EDAD 2020) y del Observatorio 2030 CSCAE--Fundación ONCE (\emph{Documenta 1.2}). Proceso actual y propuesto: comprobación manual al generar el PDF, con ajustes de maquetación. Gantt y planificación: fechas, tareas, hitos y cargas comprobados frente al original; suma total de ECTS y horas y descripciones revisadas por el autor. Metodología y objetivos: contenido, enunciados y orden de cada etapa comprobados frente a los textos del autor. Compilación con la plantilla TFUOC y revisión de la disposición por el autor.}
  \fichaprompt{Necesito diseñar un gráfico conceptual que muestre que en España, en el año 2020, según el Instituto Nacional de Estadística, había 4.38 millones de habitantes con algún tipo de diversidad funcional o discapacidad. Necesitaría que mostrases en una barra el número de personas y el porcentaje de personas con discapacidad reconocida (4.38 millones, el 9,5\% de la población). Y que ampliando el concepto, en otra barra, incluyendo aparte de las personas con discapacidad oficial, sumadas a las personas mayores de 65 años (un 7.5\% estimado) con problemas de movilidad y funcionales, personas con movilidad reducida temporal (un 2.5\%), y familias con carritos de bebé (un 10.5\%). Los datos de estos tres últimos son difíciles de estimar, pero representan un 30\% en total de la población (incluidas las personas con discapacidad reconocida). También quiero incluir que hay un porcentaje superior al 30\,\% en las personas mayores de 65 años (exactamente 32.3\%). Representa esto: primero en un gráfico conceptual los dos primeros valores, y como cifra de infografía el segundo valor}
  \fichaajuste{En lugar de dar el dato aproximado... me gustaría que mostrases todo en una columna para llegar al 30\% y así puedo poner la referencia a Barra 2 (\textgreater30\,\%): Población beneficiaria real de la accesibilidad universal * (Fuente: Observatorio 2030 CSCAE \& Fundación ONCE, Documenta 1.2)., en el gráfico conceptual}
  \fichaajuste{Lo mismo es mejor explorar la idea de hacer una infografía conl los tres valores}
  \fichaprompt{Puedes mover un poco más abajo los términos positivos de lo propuesto y hacer que entre ACTUAL manual (in situ) bien?}
  \fichaprompt{Transforma este cronograma a tu estilo por favor, dame el código latex para copiar y pegar}
  \fichaprompt{La tabla también}
  \fichaprompt{Genera un diagrama de ciclo de vida (circular), con los estilos del proyecto, para mostrar la metodología de desarrollo del proyecto: Etapas: identificación del problema (descubrir cuál es el objetivo del proyecto, basado en la lentitud de los planes de accesibilidad al tener que hacer las inspecciones in situ), diseño (análisis de las leyes, y del estado del arte a para uso de caráacterísticas de medida de accesibilidad, análisis de los datos de Albacete y Algeciras), desarrollo (implementación de los modelos y sistemas), evaluación (validacion de la precisión de los modelos, comparativa entre los modelos y las ciudades), conclusión (respuesta a las preguntas de investigación, documentación de los resultados obtenidos, líneas futuras)}
  \fichaprompt{Genera un diagrama que contenga los 6 objetivos principales del proyecto, lo que forma el MVP del mismo. Objetivos: OP1. Identificar necesidades y requisitos de accesibilidad universal, OP2. Identificar, analizar y preparar los datos, OP3. Identificar y definir las métricas e indicadores de evaluación de la accesibilidad, OP4. Entrenar modelos de aprendizaje automático de clasificación de accesibilidad, OP5. Validación de los modelos, OP6. Interpretación de los resultados. Necesito que se muestre que los objetivos se realizan en orden, según lo que te he indicado.}
\end{fichaia}

\begin{fichaia}{Figuras del marco legal}
  \fichadato{Fecha}{30/09/2026 -- 05/10/2026}
  \fichadato{Herramienta}{Claude (Anthropic)~\parencite{claude}}
  \fichadato{Fase}{WP1. Definición y objetivos, WP2.1. Análisis legal}
  \fichadato{Objetivo}{Representar visualmente los requisitos de la normativa de accesibilidad (itinerario peatonal, vados, pasos de peatones, isletas, pavimento táctil, semáforos y obras) y su evolución, para hacer menos densa la sección legal.}
  \fichadato{Material aportado}{Relación de normas redactada por el autor y tablas~\ref{tab:req-vados} y~\ref{tab:req-obras}, con los parámetros, valores y artículos seleccionados por el autor.}
  \fichadato{Resultado}{Código TikZ de las figuras~\ref{fig:ipa-plano} y~\ref{fig:ipa-seccion} (itinerario peatonal accesible, dividido en dos figuras manualmente), de la línea temporal (figura~\ref{fig:normativa_estatal}), de vados (figura~\ref{fig:vados}) y de obras (figura~\ref{fig:obras}), con cotas y leyendas.}
  \fichadato{Validación}{Itinerario peatonal: medidas contrastadas con el texto de la Orden TMA/851/2021. Línea temporal: fechas y denominaciones comprobadas en el BOE por el autor, incluidas las añadidas por la herramienta (Orden VIV/561/2010, de 1 de febrero, y publicación de la Convención en el BOE el 21 de abril de 2008). Vados y obras: valores y artículos de las cotas comprobados frente a las tablas~\ref{tab:req-vados} y~\ref{tab:req-obras} y la Orden TMA/851/2021; las dimensiones no acotadas (anchura del cruce y del vado) y la distancia entre luces son ilustrativas. Compilación con la plantilla TFUOC.}
  \fichaprompt{NEcesito que realices un mapa con medidas utilizando el componente de mapas, que defina qué es un itinerario peatona accesible según la Orden TMA/851/2021, incluyendo: ancho libre de paso (mínimo 1,80m de anchura), sin obstáculos. Altura libre de paso, mínimo 2,20m, ubicación de mobiliario en la banda exterior de la acera (con un mínimo de 40 cm de bordillo)}
  \fichaprompt{Necesito una línea temporal de la normativa estatal española:
    Ley 26/2011, de 1 de agosto: fue la ley que legisló inicialmente la norma de la ONU, aunque posteriormente fue evolucionada en las siguientes leyes y decretos.
    Real Decreto Legislativo 1/2013, de 29 de noviembre: se trata del Texto Refundido de la Ley General de derechos de las personas con discapacidad y de su inclusión social y contiene las condiciones básicas de accesibilidad y no discriminación. Es la legislación base.
    Real Decreto 505/2007, de 20 de abril: contiene el reglamento que establece las condiciones básicas de accesibilidad y no discriminación para el acceso y disfrute de los espacios públicos urbanos y edificaciones.
    Orden TMA/851/2021, de 23 de julio: es el documento técnico de referencia para este trabajo, pues establece las características técnicas y condiciones precisas de accesibilidad en las vías públicas. Entró en vigor el 2 de enero de 2022, y deroga la Orden VIV/561/2010.}
  \fichaajuste{Tras revisar la figura en la memoria, el autor pidió algunos cambios de formato: dar fondo a las etiquetas de año para que no las crucen las líneas.}
  \fichaprompt{Necesito realizar una visualización de figuras para algunos de los apartados de mi memoria, como los que te pedí para el itinerario peatonal accesible. En primer lugar, quiero un diagrama de figuras (visualización de las calles), que visualice los vados, pasos de peatones e islas [se adjunta el código de la tabla~\ref{tab:req-vados}].}
  \fichaprompt{Realiza un diagrama similar que recoja toda la normativa de obras especificado en estas normas [se adjunta el contenido de la tabla~\ref{tab:req-obras}].}
  \fichaajuste{Reducir el tamaño de las figuras, con las secciones a la derecha y la planta en el espacio restante.}
\end{fichaia}

\begin{fichaia}{Tablas y figuras de datos y modelos}
  \fichadato{Fecha}{09/10/2026 -- 17/10/2026}
  \fichadato{Herramienta}{Claude (Anthropic)~\parencite{claude}}
  \fichadato{Fase}{WP2.1. Análisis legal, WP2.2. Estado del arte}
  \fichadato{Objetivo}{Maquetar tablas y gráficos a partir de datos y textos del autor, para ahorrar trabajo de transcripción y facilitar la comparación.}
  \fichadato{Material aportado}{Resultados del cuaderno \texttt{0.03-comparison-algeciras-albacete.ipynb} y de la tabla de distribución de clases del proyecto \texttt{tfm-urban-accessibility-ml}; tablas de requisitos de la sección~\ref{sec:caracteristicas} y porcentajes de viabilidad estimados por el autor; texto del apartado~\ref{sec:ml}.}
  \fichadato{Resultado}{Código \LaTeX{} de las tablas~\ref{tab:columnas-unificadas} (columnas unificadas de Albacete y Algeciras), \ref{tab:viabilidad} (viabilidad de obtención de cada característica en imágenes aéreas y a pie de calle) y~\ref{tab:ml-revisados} (estudios de aprendizaje automático revisados), y código pgfplots de la figura~\ref{fig:distribucion-clases} (distribución de clases).}
  \fichadato{Validación}{Columnas unificadas: nombres y reglas de unificación comprobados por el autor frente al cuaderno y al conjunto de datos procesado. Viabilidad: características y porcentajes revisados frente a la sección~\ref{sec:caracteristicas}. Distribución de clases: porcentajes comprobados con la tabla obtenida en el análisis exploratorio de Albacete y Algeciras. Estudios revisados: cada celda revisada por el autor frente al texto del apartado~\ref{sec:ml}. Compilación con la plantilla TFUOC.}
  \fichaprompt{A partir del resultado de los notebook del proyecto ../tfm-urban-accessibility-ml, genera la tabla en latex con las columnas del conjunto de datos unificado de Albacete y Algeciras, indicando cómo se llama la columna unificada y cuál era la columna en el dataset original.}
  \fichaprompt{Necesito que crees una tabla para mostrar la viabilidad de las carácterísticas del apartado 2.2, poniendo una fila por cada característica "medible" para imágenes aéreas, a pie de calle. Solo pon las carácterísticas de la sección, y añade un porcentaje para aérea y a pie de calle (haz una estimación, que luego cambiaré a mano), así como una columna con la razón del porcentaje, déjame esta columna sin rellenar.}
  \fichaprompt{Crea un diagrama de barras que muestre la distribución de clases de accesibilidad de Algeciras y Albacete, siguiendo el design system, a partir de los siguientes datos [tabla con el número de tramos y el número y porcentaje de tramos de clase A, B y C en Albacete, Algeciras y el total]}
  \fichaajuste{Creo que prefiero los datos agrupados, en una sola columna (o fila)}
  \fichaprompt{Crea una tabla con label ml-revisados que liste todos los trabajos, agrupados por las subsecciones, y que tenga por columnas: modelo, métricas, y que tenga como columnas: ordinal, desbalanceo, importancia, coste y fiabilidad. Utiliza únicamente los datos que están en el apartado 3.3 para llenar la tabla.}
  \fichaajuste{Necesito que de alguna manera las columnas de ordinal, desbalanceo... no se corten}
\end{fichaia}

\begin{fichaia}{Figuras~\ref{fig:perfiles-pesos}, \ref{fig:perfiles-modelos}, \ref{fig:perfiles-multisalida} y~\ref{fig:tipos-particion}. Diagramas del estado del arte}
  \fichadato{Fecha}{18/10/2026}
  \fichadato{Herramienta}{Claude (Anthropic)~\parencite{claude}}
  \fichadato{Fase}{WP2.2 -- Estado del arte}
  \fichadato{Uso}{Generación del código TikZ de cuatro diagramas, siguiendo el sistema de diseño de la memoria: tres formas de incorporar los grupos de usuarios y tres estrategias de partición entre entrenamiento y prueba (sección~\ref{sec:transferencia-ciudades}).}
  \fichadato{Revisión}{El autor ha comprobado que cada diagrama corresponde al método descrito en los trabajos citados y a las estrategias de la sección~\ref{sec:transferencia-ciudades}, y ha ajustado su posición junto a los párrafos.}
  \fichaprompt{Crea un diagrama que muestre tres opciones diferentes (tres diagramas distintos), diferenciando entre el cálculo de pesos por grupo (Moura et al.), con indicadores de la calle como variable, en ese caso se aplican pesos distintos w\_i para cada clase de grupo. En el caso de un modelo por grupo (Dong et al.), crea un diagrama que muestre que para cada grupo hay un modelo (M\_i), y la salida del modelo es una clase de grupo, y finalmente otro de un modelo con varias salidas por grupo, quiero que indiques que a partir de las variables se asignan varias salidas por grupo, en este caso continuas. Necesito que el diagrama no ocupe demasiado, para poder mostrarse junto con los párrafos.}
  \fichaprompt{Crea un diagrama que muestre cómo se dividen los tramos de calles en: partición aleatoria agrupada por calle o barrio o entre ciudades, en los conjuntos de entrenamiento o test, que tenga un tamaño pequeño para ponerlo al lado de un párrafo y que sea más visual). En partición aleatoria, los tramos de una misma calle pueden caer en test o entrenamiento; en la agrupada todos los tramos de una calle caen en entrenamiento o test, pero nunca separadas; mientras que entre ciudades, se entrenan todas las calles en una ciudad, y se prueban en otra. Puedes poner cuadritos con dos colores (entrenamiento y prueba) o algo similar. Sigue el design system.}
\end{fichaia}

\begin{fichaia}{Figura~\ref{fig:ponderacion-submuestreo}. Ponderación de clases frente a submuestreo}
  \fichadato{Fecha}{23/10/2026}
  \fichadato{Herramienta}{Claude (Anthropic)~\parencite{claude}}
  \fichadato{Objetivo}{Representar gráficamente la decisión de tratamiento del desbalanceo de la subsección~\ref{subsec:ml-ordinal}, en la parte más densa del estado del arte.}
  \fichadato{Material aportado}{Distribución de las clases de accesibilidad de Albacete (figura~\ref{fig:distribucion-clases}) y sistema de diseño de la memoria (\texttt{tfm-estilo.sty}).}
  \fichaprompt{Necesito un gráfico usando el sistema de diseño de la memoria, que compare la ponderación de clases con el submuestreo usando los datos de Albacete, que use los datos reales de Albacete del anterior gráfico.}
  \fichaajuste{Cambia los pesos numéricos por valores relativos (alto, medio y bajo).}
\end{fichaia}

\begin{fichaia}{Figuras~\ref{fig:arquitectura-ingesta} y~\ref{fig:arquitectura-modelado}. Diagrama de la arquitectura del sistema}
  \fichadato{Fecha}{25/10/2026}
  \fichadato{Herramienta}{Claude (Anthropic)~\parencite{claude}}
  \fichadato{Fase}{WP2 -- Redacción de la memoria}
  \fichaprompt{Transform the diagram DiagramaArquitectura.drawio.png to make it look like the system design of the project}
  \fichaajuste{Redibujar en TikZ como diagrama de flujo, dividirlo en dos figuras y compactarlo en vertical, y corregir detalles de formato (rótulo superpuesto, informe en lugar de base de datos).}
\end{fichaia}

\begin{fichaia}{Consultas y análisis de la literatura}
  \fichadato{Fecha}{Octubre de 2026 (varias sesiones)}
  \fichadato{Herramienta}{Gemini (Google)~\parencite{gemini} y Claude (Anthropic)~\parencite{claude}}
  \fichadato{Fase}{WP2 -- Contexto y estado del arte}
  \fichadato{Uso}{Propuesta de cadenas de búsqueda para Web of Science y Google Scholar (tabla~\ref{tab:cadenas-busqueda}); análisis de la literatura seleccionada (algoritmos multiobjetivo, métricas y estrategias de partición espacial y transferencia entre ciudades); y resolución de dudas puntuales sobre las fuentes de datos, el preprocesado y la validación.}
  \fichadato{Revisión}{El autor ejecutó las cadenas en ambas fuentes, descartó las que devolvían resultados poco relevantes y seleccionó los estudios según los criterios de la sección~\ref{subsec:criterios-seleccion}.}
  \fichaprompt{Recomiéndame palabras y frases de búsqueda, varias diferentes por cada uno de los análisis a realizar, que sea un listado de cadenas de al menos 5 por cada uno.}
  \fichaprompt{¿Puedes resumir las 3 fases del algoritmo MTBR de Lin et al. (2019) y cómo se compara con MOTC?}
  \fichaprompt{¿Qué métricas se utilizan en las fuentes seleccionadas para evaluar modelos de \emph{machine learning} multiobjetivo y la evaluación de la caminabilidad?}
  \fichaprompt{Identifica qué fuentes evalúan modelos entre distintas ciudades o zonas geográficas, indicando la estrategia espacial, las métricas y los resultados.}
  \fichaprompt{¿Cuáles fueron los conjuntos de datos, la estrategia de generalización entre ciudades y las métricas utilizadas en el estudio de Weld et al. (2019)?}
  \fichaprompt{En el estudio de Wu et al. (2026), ¿qué estrategia se sigue y cuál fue el resultado de la validación agrupada por calle frente a la división aleatoria?}
  \fichaprompt{¿Qué estrategia sigue Li et al. (2024) para escalar la predicción de la caminabilidad desde un distrito muestra hasta la escala de toda la ciudad?}
  \fichaprompt{¿Qué ocurre con el rendimiento de un modelo cuando se entrena en una ciudad y se prueba en otra con infraestructura urbana diferente (por ejemplo, en Duan et al., 2022)?}
  \fichaprompt{¿Crees que es posible que los datos de pendiente transversal se puedan obtener por fuentes cartográficas públicas?}
  \fichaprompt{¿Qué diferencia hay entre una foto aérea y una foto satelital?}
  \fichaprompt{¿Qué me puedo descargar de aquí? https://pnoa.ign.es/web/portal/pnoa-lidar/productos-a-descarga}
  \fichaprompt{Pensaba que los parámetros de entrada de un modelo tenían que quedarse siempre en la misma escala. ¿Es así?}
  \fichaprompt{¿Qué importancia tiene en mi trabajo la fiabilidad? ¿Quizás cuando no se puedan tener datos por la vegetación?}
  \fichaprompt{¿Y si busco un mapa de barriadas de Algeciras, para poder identificar las calles separadas por barriada?}
\end{fichaia}

\begin{fichaia}{Bibliografía de la memoria}
  \fichadato{Fecha}{30/09/2026 -- Octubre de 2026}
  \fichadato{Herramienta}{Gemini (Google)~\parencite{gemini} y Claude (Anthropic)~\parencite{claude}}
  \fichadato{Fase}{WP1. Definición y objetivos, WP2 -- Redacción de la memoria}
  \fichadato{Objetivo}{Convertir la bibliografía manual (\texttt{thebibliography}) a BibTeX y validar las referencias y el formato de las citas.}
  \fichadato{Resultado}{Archivo \texttt{referencias.bib}, configuración de \texttt{biblatex} en el preámbulo y \texttt{\textbackslash printbibliography[heading=bibintoc]}; correcciones de entradas y de citas secundarias.}
  \fichadato{Validación}{Compilación con biber; verificación de citas en el documento y de que todas las referencias se generan correctamente en el PDF; revisión por el autor de las entradas y del año de publicación de cada una.}
  \fichaprompt{Transforma la bibliografía en referencias.bib, y carga la bibliografía en la sección en la que está}
  \fichaprompt{Comprueba si estas entradas BibTeX para mastelini2019 y lin2019 son correctas y contienen todos los campos de la publicación.}
  \fichaprompt{Revisa si esta entrada BibTeX de Duan et al. (2022) en ACM SIGACCESS contiene la sintaxis correcta para actas de congreso (\texttt{booktitle}, \texttt{doi}, \texttt{pages}).}
  \fichaprompt{Comprueba si el año de publicación asignado al artículo de Lin et al. en TKDE corresponde a la versión preliminar (2019) o al volumen definitivo de la revista (2021).}
\end{fichaia}
